% Exercises the citation commands beyond \cite: biblatex's \textcite,
% \parencite, \autocite (and a starred form), \citeauthor as a representative of
% the commands printing only a fragment of a citation, and a command of one's
% own registered with \proofgraphcitecommand. The last proof cites one work
% three ways at once, checking that this still yields a single node and edge.
% citelabel=key keeps the reference independent of the bibliography style.
\documentclass{article}
\usepackage{amsthm}
\usepackage[backend=biber,style=numeric]{biblatex}
\addbibresource{refs.bib}
\usepackage[cite,citelabel=key]{proofgraph}
\newcommand\mycite[1]{(see~#1)}
\proofgraphcitecommand{mycite}
\newtheorem{theorem}{Theorem}
\begin{document}
\begin{theorem}\label{t1}One.\end{theorem}
\begin{proof}\textcite{knuth} and \parencite[Thm~3]{turing}.\end{proof}
\begin{theorem}\label{t2}Two.\end{theorem}
\begin{proof}\autocite{euclid} and \textcite*{knuth}.\end{proof}
\begin{theorem}\label{t3}Three.\end{theorem}
\begin{proof}\citeauthor{turing} and \parencite[see][p.~7]{euclid}.\end{proof}
\begin{theorem}\label{t4}Four.\end{theorem}
\begin{proof}\mycite{euclid}; \citeauthor{knuth} in \textcite{knuth}.\end{proof}
\printbibliography
\end{document}
